\section{Diego Alberto Vázquez Soqui}
Comprender los fundamentos teóricos de la robótica es clave para desarrollar sistemas que funcionen de manera eficiente y precisa. La cinemática permite analizar y diseñar el movimiento de los robots, mientras que la dinámica ofrece una visión más completa al considerar las fuerzas involucradas. El control, por su parte, garantiza que el robot ejecute sus movimientos de forma estable y confiable. Además, herramientas como ROS han transformado el desarrollo robótico al proporcionar una estructura flexible y escalable para integrar sensores, actuadores y algoritmos. En conjunto, estos conceptos forman la base para abordar desafíos reales en la automatización y la inteligencia artificial aplicada a la robótica.